\documentclass[10pt,a4paper]{amsart}
\usepackage[margin=19mm]{geometry}
\usepackage{amsmath,amssymb,mathtools}
\usepackage[scr=rsfs,cal=euler]{mathalfa}
\usepackage{xfrac}
\usepackage[unicode,hidelinks,bookmarks=false]{hyperref}
\newcommand{\ie}{i.e.\ }
\newcommand{\cf}{cf.\ }
\newcommand{\resp}{respectively\ }
\newcommand{\Subsection}{\subsection}
\providecommand{\nobreakdash}{\nobreak\hbox{-}}
\setlength{\emergencystretch}{2em}
\multlinegap0pt
\newtheorem{theorem}{Theorem}
\title{Existence and maximal corank of simple $Z_p$-invariant germs}
\author{Ivan Proskurnin}
\date{Source: arXiv:2604.26659v2 (CC BY 4.0)}
\begin{document}
\maketitle

\noindent\emph{Source and licence.} Existence and maximal corank of simple $Z_p$-invariant germs. Version arXiv:2604.26659v2 (CC BY 4.0), distributed by arXiv under the Creative Commons Attribution 4.0 International licence, http://creativecommons.org/licenses/by/4.0/, as recorded in the arXiv licence metadata for this exact version and shown on the abstract page https://arxiv.org/abs/2604.26659v2. Full licence notice: \texttt{licenses/arxiv-2604.26659-CC-BY-4.0.txt}.\par\smallskip

\noindent\textit{Editorial scope.} Counted unit: the article's theorem theorem (source0.tex lines 159--161). The article's complete original proof of it is delivered with it (source0.tex lines 163--179, one \texttt{proof} environment of 17 lines). No mathematics is written, repaired or re-derived: the statement and the proof are the article's own lines, in the article's own notation, and no retained line is substituted or re-flowed.  No further source-local context is needed: every cross-reference of every retained line resolves inside the two delivered spans.  Nothing was omitted from any retained span, so the package carries no omission record.  The retained lines cite 1 works, whose entries are transcribed verbatim from the article's own bibliography source in the e-print and delivered with short editorial labels recorded in the item's reference mapping.  No retained line contains a figure environment, an image-inclusion command or drawing code, so no image is delivered and the package declares no asset dependency.  Editorial formatting only, in addition: an AMS \texttt{amsart} preamble that restates the article's own declarations and macros used by the retained lines, namely the environments \texttt{theorem} on separate counters, together with the standard packages those lines need.  The retained environments print in this document as Theorem 1 in order of appearance, whereas the article prints the same items with the numbering of its own full document.  The complete native source of the article as distributed in the arXiv e-print for this exact version is bundled unchanged as \texttt{sources/199409/source0.tex}.
\begin{theorem}
$f$ is an equivariantly stable $G_n$-invariant germ.
\end{theorem}

\begin{proof}

Consider the case of odd $n$, so that $G_n= \mathbb{Z}_{d_1\ldots d_n + 1}$. To prove that $f$ is stable we need to establish that $\nu(f)=1$, i.e. that the only invariant element in some monomial basis of $Q_f$ is $1$. The monomial basis of $Q_f$ consists (see, e.g. \cite{10}) of monomials $x_1^{r_1}\ldots x_n^{r_n}, 0 \leq r_i < d_i$. Proving that such monomial is only invariant when $r_1 = r_2 = \ldots = r_n =0$ is obviously equivalend to proving that 

$r_1 - r_2 d_1 + \ldots + (-1)^{n-1} r_n d_1 d_2 \ldots d_{n-1} \equiv 0 (mod \; d_1\ldots d_n + 1)$ only if $r_1 = r_2 = \ldots = r_n =0$ for $r_i < d_i$.


 $| r_1 - r_2 d_1 + \ldots + (-1)^{n-1}  r_n d_1 d_2 \ldots d_{n-1} | \leq | r_1| + |r_2 d_1| + \ldots + | r_n d_1 d_2 \ldots d_{n-1} |$. As $r_i < d_i$  $| r_1| + |r_2 d_1| + \ldots + | r_n d_1 d_2 \ldots d_{n-1} | \leq (d_1-1) + (d_2-1)d_1 + \ldots + (d_n - 1) d_1 d_2 \ldots d_{n-1}$. The last sum is a telescopic series equal to $d_1 d_2 \ldots d_n - 1$. 
 
 Therefore $r_1 - r_2 d_1 + \ldots + (-1)^{n-1}  r_n d_1 d_2 \ldots d_{n-1} \equiv 0 (mod \; d_1\ldots d_n + 1)$ 
 
 only if $r_1 - r_2 d_1 + \ldots + (-1)^{n-1}  r_n d_1 d_2 \ldots d_{n-1} = 0$. 
 
 By a similar argument, the sum of first $n-1$ terms in the equality $r_1 - r_2 d_1 + \ldots + (-1)^{n-1}  r_n d_1 d_2 \ldots d_{n-1} = 0$ has absolute value at most $d_1 d_2 \ldots d_{n-1} - 1$, so $r_1 - r_2 d_1 + \ldots + (-1)^{n-1}  r_n d_1 d_2 \ldots d_{n-1} = 0$ means $r_n=0$. Proceeding inductively, we obtain $r_1 = r_2 = \ldots = r_n =0$. 
 
 The case of even $n$ is basically identical, it only needs to be noted that the inequality  $| r_1 - r_2 d_1 + \ldots + (-1)^{n-1}  r_n d_1 d_2 \ldots d_{n-1} | \leq d_1 \ldots d_n -1$ is actually strict, as the sum $r_1 - r_2 d_1 + \ldots + (-1)^{n-1}  r_n d_1 d_2 \ldots d_{n-1}$ always has either some terms of different signs or some zero terms.
\end{proof}

\begin{thebibliography}{99}
\bibitem[10]{10} A.~Basalaev, ``Mirror symmetry for invertible singularities'', \url{https://www.hse.ru/mirror/pubs/share/1062447379.pdf}.
\end{thebibliography}
\end{document}
